\section*{Preface to Volume I}
\addcontentsline{toc}{section}{Preface to Volume I}

Why another book on quantum field theory? Today the student of quantum field
theory can choose from among a score of excellent books, several of them quite
up-to-date. Another book will be worth while only if it offers something new in
content or perspective.

As to content, although this book contains a good amount of new material, I
suppose the most distinctive thing about it is its generality; I have tried
throughout to discuss matters in a context that is as general as possible. This
is in part because quantum field theory has found applications far removed from
the scene of its old successes, quantum electrodynamics, but even more because I
think that this generality will help to keep the important points from being
submerged in the technicalities of specific theories. Of course, specific
examples are frequently used to illustrate general points, examples that are
chosen from contemporary particle physics or nuclear physics as well as from
quantum electrodynamics.

It is, however, the perspective of this book, rather than its content, that
provided my chief motivation in writing it. I aim to present quantum field
theory in a manner that will give the reader the clearest possible idea of
\emph{why} this theory takes the form it does, and why in this form it does such
a good job of describing the real world.

The traditional approach, since the first papers of Heisenberg and Pauli on
general quantum field theory, has been to take the existence of fields for
granted, relying for justification on our experience with electromagnetism, and
`quantize' them --- that is, apply to various simple field theories the rules of
canonical quantization or path integration. Some of this traditional approach
will be found here in the historical introduction presented in Chapter 1. This
is certainly a way of getting rapidly into the subject, but it seems to me that
it leaves the reflective reader with too many unanswered questions. Why should
we believe in the rules of canonical quantization or path integration? Why
should we adopt the simple field equations and Lagrangians that are found in the
literature? For that matter, why have fields at all? It does not seem
satisfactory to me to appeal to experience; after all, our purpose in
theoretical physics is not just to describe the world as we find it, but to
explain --- in terms of a few fundamental principles --- why the world is the
way it is.

The point of view of this book is that quantum field theory is the way it is
because (aside from theories like string theory that have an infinite number of
particle types) it is the only way to reconcile the principles of quantum
mechanics (including the cluster decomposition property) with those of special
relativity. This is a point of view I have held for many years, but it is also
one that has become newly appropriate. We have learned in recent years to think
of our successful quantum field theories, including quantum electrodynamics, as
`effective field theories,' low-energy approximations to a deeper theory that
may not even be a field theory, but something different like a string theory. On
this basis, the reason that quantum field theories describe physics at
accessible energies is that \emph{any} relativistic quantum theory will look at
sufficiently low energy like a quantum field theory. It is therefore important
to understand the rationale for quantum field theory in terms of the principles
of relativity and quantum mechanics. Also, we think differently now about some
of the problems of quantum field theories, such as non-renormalizability and
`triviality,' that used to bother us when we thought of these theories as truly
fundamental, and the discussions here will reflect these changes. This is
intended to be a book on quantum field theory for the era of effective field
theories.

The most immediate and certain consequences of relativity and quantum mechanics
are the properties of particle states, so here particles come first --- they
are introduced in Chapter 2 as ingredients in the representation of the
Poincar\'e group in the Hilbert space of quantum mechanics. Chapter
3 provides a framework for addressing the fundamental dynamical question:
given a state that in the distant past looks like a certain collection of free
particles, what will it look like in the future? Knowing the generator of
time-translations, the Hamiltonian, we can answer this question through the
perturbative expansion for the array of transition amplitudes known as the
$S$-matrix. In Chapter 4 the principle of cluster decomposition is invoked to
describe how the generator of time-translations, the Hamiltonian, is to be
constructed from creation and annihilation operators. Then in Chapter 5 we
return to Lorentz invariance, and show that it requires these creation and
annihilation operators to be grouped together in causal quantum fields. As a
spin-off, we deduce the CPT theorem and the connection between spin and
statistics. The formalism is used in Chapter 6 to derive the Feynman rules for
calculating the $S$-matrix.

It is not until Chapter 7 that we come to Lagrangians and the canonical
formalism. The rationale here for introducing them is not that they have proved
useful elsewhere in physics (never a very satisfying explanation) but rather
that this formalism makes it easy to choose interaction Hamiltonians for which
the $S$-matrix satisfies various assumed symmetries. In particular, the Lorentz
invariance of the Lagrangian density ensures the existence of a set of ten
operators that satisfy the algebra of the Poincar\'e group and, as we show in
Chapter 3, this is the key condition that we need to prove the Lorentz
invariance of the $S$-matrix. Quantum electrodynamics finally appears in Chapter
8. Path integration is introduced in Chapter 9, and used to justify some of the
hand-waving in Chapter 8 regarding the Feynman rules for quantum
electrodynamics. This is a somewhat later introduction of path integrals than is
fashionable these days, but it seems to me that although path integration is by
far the best way of rapidly deriving Feynman rules from a given Lagrangian, it
rather obscures the quantum mechanical reasons underlying these calculations.

Volume I concludes with a series of chapters, 10--14, that provide an
introduction to the calculation of radiative corrections, involving loop
graphs, in general field theories. Here too the arrangement is a bit unusual; we
start with a chapter on non-perturbative methods, in part because the results we
obtain help us to understand the necessity for field and mass renormalization,
without regard to whether the theory contains infinities or not. Chapter 11
presents the classic one-loop calculations of quantum chromodynamics, both as an
opportunity to explain useful calculational techniques (Feynman parameters,
Wick rotation, dimensional and Pauli--Villars regularization), and also as a
concrete example of renormalization in action. The experience gained in Chapter
11 is extended to all orders and general theories in Chapter 12, which also
describes the modern view of non-renormalizability that is appropriate to
effective field theories. Chapter 13 is a digression on the special problems
raised by massless particles of low energy or parallel momenta. The Dirac
equation for an electron in an external electromagnetic field, which
historically appeared almost at the very start of relativistic quantum
mechanics, is not seen here until Chapter 14, on bound state problems, because
this equation should not be viewed (as Dirac did) as a relativistic version of
the Schr\"odinger equation, but rather as an approximation to a true relativistic
quantum theory, the quantum field theory of photons and electrons. This chapter
ends with a treatment of the Lamb shift, bringing the confrontation of theory
and experiment up to date.

The reader may feel that some of the topics treated here, especially in Chapter
3, could more properly have been left to textbooks on nuclear or elementary
particle physics. So they might, but in my experience these topics are usually
either not covered or covered poorly, using specific dynamical models rather
than the general principles of symmetry and quantum mechanics. I have met
string theorists who have never heard of the relation between time-reversal
invariance and final-state phase shifts, and nuclear theorists who do not
understand why resonances are governed by the Breit--Wigner formula. So in the
early chapters I have tried to err on the side of inclusion rather than
exclusion.

Volume II will deal with the advances that have revived quantum field theory in
recent years: non-Abelian gauge theories, the renormalization group, broken
symmetries, anomalies, instantons, and so on.

I have tried to give citations both to the classic papers in the quantum theory
of fields and to useful references on topics that are mentioned but not
presented in detail in this book. I did not always know who was responsible for
material presented here, and the mere absence of a citation should not be taken
as a claim that the material presented here is original. But some of it is. I
hope that I have improved on the original literature or standard textbook
treatments in several places, as for instance in the proof that symmetry
operators are either unitary or antiunitary; the discussion of superselection
rules; the analysis of particle degeneracy associated with unconventional
representations of inversions; the use of the cluster decomposition principle;
the derivation of the reduction formula; the derivation of the external field
approximation; and even the calculation of the Lamb shift.

I have also supplied problems for each chapter except the first. Some of these
problems aim simply at providing exercise in the use of techniques described in
the chapter; others are intended to suggest extensions of the results of the
chapter to a wider class of theories.

In teaching quantum field theory, I have found that each of the two volumes of
this book provides enough material for a one-year course for graduate students.
I intended that this book should be accessible to students who are familiar
with non-relativistic quantum mechanics and classical electrodynamics. I assume
a basic knowledge of complex analysis and matrix algebra, but topics in group
theory and topology are explained where they are introduced.

This is not a book for the student who wants immediately to begin calculating
Feynman graphs in the standard model of weak, electromagnetic, and strong
interactions. Nor is this a book for those who seek a higher level of
mathematical rigor. Indeed, there are parts of this book whose lack of rigor
will bring tears to the eyes of the mathematically inclined reader. Rather, I
hope it will suit the physicists and physics students who want to understand
\emph{why} quantum field theory is the way it is, so that they will be ready for
whatever new developments in physics may take us beyond our present
understandings.

\begin{center}
  *\quad*\quad*
\end{center}

Much of the material in this book I learned from my interactions over the years
with numerous other physicists, far too many to name here. But I must
acknowledge my special intellectual debt to Sidney Coleman, and to my colleagues
at the University of Texas: Arno Bohm, Luis Boya, Phil Candelas, Bryce DeWitt,
Cecile DeWitt-Morette, Jacques Distler, Willy Fischler, Josh Feinberg, Joaquim
Gomis, Vadim Kaplunovsky, Joe Polchinski, and Paul Shapiro. I owe thanks for
help in the preparation of the historical introduction to Gerry Holton, Arthur
Miller, and Sam Schweber. Thanks are also due to Alyce Wilson, who prepared the
illustrations and typed the \LaTeX{} input files until I learned how to do it,
and to Terry Riley for finding countless books and articles. I am grateful to
Maureen Storey and Alison Woollatt of Cambridge University Press for helping to
ready this book for publication, and especially to my editor, Rufus Neal, for
his friendly good advice.

\vspace{2.5em}
\noindent
\begin{minipage}[t]{0.45\textwidth}
Austin, Texas\\
October, 1994
\end{minipage}
\hfill
\begin{minipage}[t]{0.45\textwidth}
\raggedleft
\textsc{Steven Weinberg}
\end{minipage}

